%%%%%%%%%%%%%%%%%%%%%%%%%%%%%%%%%%%%%%%%%%%%%%%%%%
%%%%%%%%%%%%%%%%%%%%%%%%%%%%%%%%%%%%%%%%%%%%%%%%%%
\section{Introduction}
\label{remdo:sec:intro}

%%%%%%%%%%%%%%%%%%%%%%%%%%%%%%%%%%%%%%%%%%%%%%%%%%

%(Why WECs, in renewable energy context):

% stronger opening sentence
\ifdefined\DISSERTATION
    Reducing global greenhouse emissions requires every viable carbon-free energy source, including ocean wave energy.
    Wave power offers higher temporal consistency and energy density than wind or solar, and its complementarity with other renewables can ease energy storage requirements and improve capacity adequacy \citep{akdemir_opportunities_2023,bhattacharya_timing_2021,pennock_temporal_2022}.
    Wave energy converters (WECs) can deliver both utility-scale electricity and power for offshore applications such as aquaculture, desalination, and marine sensing \citep{livecchi_powering_2019}.
    Yet WECs still face a substantial cost gap relative to wind and solar before large-scale deployment becomes viable, motivating a design process focused on techno-economic competitiveness from the earliest stages.

    Achieving that competitiveness is complicated by strong interdisciplinary coupling.
    The powertrain, controller, bulk dimensions, hydrodynamic shape, and structural thicknesses must be selected jointly to balance power, cost, and survivability while remaining practically feasible.
    Sequential discipline-by-discipline design cannot exploit the coupling between subsystems, and computational cost has historically prevented system-level concurrent optimization from filling that gap.
    Automating concurrent decision-making allows WEC designers to systematically explore tradeoffs across major design choices and, ultimately, to compare architectures on a standardized footing.

    A landmark step toward a multidisciplinary WEC design baseline was the U.S.\ Department of Energy's Reference Model Project (2010--2014), which produced six benchmark designs along with design artifacts and a comprehensive report \citep{RM3}.\footnote{\url{https://openei.org/wiki/PRIMRE/Signature_Projects/Reference_Model}}
    The third reference model (RM3), a two-body point-absorber WEC sited for Humboldt Bay, California, has received considerable research attention.
    Its surface float oscillates against a spar consisting of a vertical column and a subsurface damping plate (\Cref{remdo:fig:rm3-parts}), with electrical power produced from the relative heave motion.
    The present study builds on that work, presenting a multidisciplinary techno-economic optimization of RM3.
\else
    Wave energy converters (WECs) can deliver both utility-scale electricity and power for offshore applications such as aquaculture, desalination, and marine sensing \citep{livecchi_powering_2019,akdemir_opportunities_2023,bhattacharya_timing_2021,pennock_temporal_2022}, but large-scale deployment requires substantial cost reduction.
    Achieving that reduction is complicated by strong interdisciplinary coupling: the powertrain, controller, bulk dimensions, hydrodynamic shape, and structural thicknesses jointly determine performance, cost, and survivability.
    Sequential single-discipline design cannot exploit this coupling, motivating system-level concurrent optimization.

    This paper presents a multidisciplinary techno-economic optimization of the Reference Model 3 (RM3) WEC \citep{RM3}: a two-body point absorber sited for Humboldt Bay, California, in which a surface float oscillates against a spar with a subsurface damping plate (\Cref{remdo:fig:rm3-parts}), with electrical power produced from the relative heave motion.
\fi

\begin{figure}[!ht]
\centering
\includegraphics[\FigureOptions{re:figs/from-matlab/RM3_image.pdf}]{figs/from-matlab/RM3_image.pdf}
\caption{RM3 geometry visualized}
\label{remdo:fig:rm3-parts}
\end{figure}

The remainder of the introduction reviews WEC optimization literature, frames the research gap, articulates the case for multidisciplinary design optimization, and identifies the paper's contributions and structure.

%%%%%%%%%%%%%%%%%%%%%%%%%%%%%%%%%%%%%%%%%%%%%%%%%%
\subsection{WEC Design Optimization: State of the Art and Gaps}
\label{remdo:sec:lit}

\ifdefined\DISSERTATION
    Before reviewing the academic state of the art, it is worth assessing how optimization and concurrent design are practiced in the wave energy industry.
    \citet{trueworthy_wave_2020} surveys 25 WEC designers and developers to report that industry use of optimization is far from universal, with the highest adoption rates at 48\% and 40\% for controls and hydrodynamic shape respectively.
    The authors attribute this to the time, expertise, and computational resources optimization demands; the inability of optimization to capture all WEC requirements; and the diversity of available techniques.
    Half of respondents design all subsystems concurrently; the rest follow a sequential pattern (shape and PTO first, then controller and moorings, then power transportation).
    The same study concludes that concurrent design is promising but under-utilized, and emphasizes the importance of incorporating all design requirements at an early stage.
    
    Several review papers describe the state of the art in WEC design optimization.
    \Cref{remdo:tab:lit} draws on these and the underlying primary literature \citep{garcia-teruel_review_2021,guo_geometric_2021,ringwood_empowering_2023,clark_reliability-based_2018,coe_survey_2018,ove_arup__partners_ltd_structural_2016,paduano_towards_2024,giannini_wave_2022} to compare the most relevant studies along modeling and optimization axes.
\else
    After surveying 25 WEC designers, \citet{trueworthy_wave_2020} reports that only 48\% and 40\% of developers use optimization for controls and hydrodynamic shape respectively, and only 50\% design all subsystems concurrently due to the time, expertise, and computational resources that optimization demands.
    They conclude that concurrent design is a promising but under-utilized technique.
    Review papers \citep{garcia-teruel_review_2021,guo_geometric_2021,ringwood_empowering_2023,clark_reliability-based_2018} and primary literature on WEC survivability guidelines \citep{coe_survey_2018,ove_arup__partners_ltd_structural_2016,paduano_towards_2024,giannini_wave_2022} together motivate the comparison in \Cref{remdo:tab:lit}.
\fi

\Cref{remdo:tab:lit} compares previous WEC optimization studies along modeling fidelity (hydrodynamics, drag, degrees of freedom, domain, controls, mooring, powertrain, structures, economics, sea states) and optimization formulation (optimizer, objectives, constraints, design variables).
Column-abbreviation legends are provided in \Cref{remdo:sec:appendix-lit-legend}.

\ifdefined\DISSERTATION
    Point absorbers dominate as the most common architecture, though all major device types are represented and several studies optimize arrays.
    Boundary element method (BEM) hydrodynamics is most common, with occasional semi-analytical (MEEM) treatments and rare CFD use.
    Dynamic modeling occurs in either frequency (F) or time (T) domain, with constraint-handling extensions (F+) and pseudo-spectral (PS) methods both gaining traction.
    Controllers range from pure damping (P) and proportional-integral (PI) through their constrained extensions (P+/PI+) to nonlinear structured (NL) and unstructured (UNS) forms.
    Mooring, powertrain, structural, and economic disciplines are unevenly represented.
    Optimizers used include genetic algorithms (GA), local optimizers (LOC), Pontryagin's maximum principle (PMP), design of experiments (DOE), manual iteration (IT), and brute force sweeps (BF), with objectives spanning power (PWR), economic viability (ECON), variability (VAR), array interaction (ARRY), load (LO), and bandwidth (BW).
    The constraints column captures all enforced requirements whether implemented as optimization constraints or as model assumptions, spanning geometry (GEO), motion amplitude (AMP), PTO (PTO), stability (STA), structural failure (STR), and power (PWR).
    The design-variable count includes only true design parameters, not state variables internal to direct-transcription or pseudo-spectral formulations.
    Bulk dimensions (DIM), hull shape (SHP), and control parameters (CTL) are common; structural (STR) variables are rare; and most studies have fewer than five design variables, with the three largest at 14, 15, and 66.
\else
    Point absorbers dominate the literature, and BEM hydrodynamics is the most commonly used method, with occasional semi-analytical (MEEM) approaches.
    Dynamic modeling is typically performed in either the frequency (F) or time domain (T), sometimes with constraint-handling extensions (F+) or pseudo-spectral (PS) variants.
    Mooring, powertrain, structural, and economic disciplines are included inconsistently across studies.
    Common optimization methods include genetic algorithms (GA), local optimizers (LOC), Pontryagin’s maximum principle (PMP), and design-of-experiment (DOE) or brute-force (BF) approaches.
    Most studies consider fewer than five design variables, although the three largest cases include 14, 15, and 66 variables.
\fi
%Need to copy sources over from slide 27 and 35 of F24 reseach slides with original charts (with color)

%\begin{landscape}
\begingroup
\centering

%\begin{tabular}
\singleColMacro{
\setlength{\tabcolsep}{1pt}
\expanded{
    \noexpand\longtable{\TableColumnSpecs{re:lit-review}}
}
\rot{\textbf{Ref}} & \rot{\textbf{Device}} & \rot{\textbf{Hydro}}& \rot{\textbf{Drag}} & \rot{\textbf{\# DOF}} & \rot{\textbf{Domain}} & \rot{\textbf{Controls}} & \rot{\textbf{Mooring}} & \rot{\textbf{Powertrain}} & \rot{\textbf{Structures}} & \rot{\textbf{Economics}} & \rot{\textbf{Sea state}}
& \rot{\textbf{Optimizer}}& \rot{\textbf{Objective}} & \rot{\textbf{Constraints}} & \rot{\textbf{\shortstack{Design \\ Variables}}}& \rot{\textbf{\# DVs}} \\
\hline\hline
This work & 2-PA & MEEM & DF & 1 & F+ & PI+ & - & DYN & AN & STR, PTO & REG+, STO
& LOC & ECON, PWR & GEO, AMP, PTO, STA, STR & DIM, STR, PTO & 12 \\

\citep{mccabe_multidisciplinary_2022} & 2-PA & ALG & - & 1 & F+ & PI+ & - & EFF & AN & STR & REG+, STO
& LOC & ECON, VAR & GEO, PTO & DIM, PTO & 7 \\

\citep{khanal_multi-objective_2024} & 1-PA array & BEM & - & 4 & F & PI+ & - & - & - & GEO & REG
& GA & ECON, ARRY & GEO, PTO & DIM, CTL & 14 \\

\citep{gaudin_single_2021} & 1-PA array & MEEM & - & 120 & F & PI & DES & - & AN & MOOR & IRR+, STO & GA & PWR, ECON & GEO, STR & DIM, CTL & 4 \\

\citep{edwards_optimisation_2022} & 1-PA & BEM & - & 1 & F & P & - & - & - & GEO & REG
& GA & ECON & AMP, GEO, PWR, STA & DIM, SHP & 15 \\

\citep{garcia-teruel_reliability-based_2021}& 1-PA & BEM & LD & 1 & F+ & PI & - & - & LD & - & IRR+
& GA & PWR, LD & AMP, PTO & DIM, SHP & 66 \\

\citep{garcia-teruel_design_2022}& 1-PA & BEM & - & 2 & F+ & PI & - & - & - & - & IRR
& GA & ECON, PWR & AMP, PTO & DIM, SHP & 4 \\

\citep{cotten_multi-objective_2022} & ATN & BEM & - & 10 & F+ & P & - & - & LD & - & IRR+
& GA & ECON, LD & GEO, STA & DIM & 8 \\

\citep{abdulkadir_control_2024} & 1-PA & MEEM & LD & 2 & T & NL P & - & - & - & - & IRR
& GA & PWR & GEO & DIM, SHP, CTL & 8 \\

\citep{housner_numerical_2024} & TRM & BEM & - & 4 & T & PI & DYN & - & - & - & REG
& LOC & PWR & GEO & DIM, SHP, CTL & 3 \\

\citep{al_shami_parameter_2019} & 2-PA & BEM & DF & 2 & F & PI & - & - & - & - & REG
& DOE & PWR, BW & - & DIM, CTL, SHP & 7 \\

\citep{RM3} & 2-PA & BEM & Q & 2 & T & P & DES & - & FEA & FULL & IRR+, STO
& IT & - & - & - & - \\

\citep{mi_multi-scale_2025} & OSW & BEM & Var. & 8 & T & PI & DYN, DES & DYN & FEA & FULL & IRR+, STO
& IT & - & - & - & - \\

\citep{an_optimal_2024} & OTP & CFD & CFD & - & T & - & - & - & FEA & - & STO
& ALT & ECON & STR & STR & 4 \\

\citep{ambuhl_reliability-based_2014} & 1-PA & - & Q & - & - & - & DES & - & AN & STR & STO
& BF & ECON & STR, STA & DIM, STR & 3 \\

\citep{nguyen_theoretical_2024} & OSW & MEEM & LD & 1 & F & PI & - & - & LD & - & IRR
& BF & ECON, LD & - & DIM & 2 \\

\citep{ferri_balancing_2014} & 1-PA & BEM & LD & 1 & T & P, PI+, UNS & - & DYN & LD & STR & IRR+
& LOC & PWR, LD & AMP, PTO & CTL & 3 \\

\citep{rosati_control_2023} & OWC & BEM & - & 1 & T & NL P+ & - & EFF & - & PTO & IRR+
& BF & ECON, VAR & PTO & PTO, CTL & 2 \\

\citep{son_performance_2016} & 2-PA & MEEM & LD & 1 & F & P & - & EFF & - & - & REG
& BF & PWR & - & PTO & 2 \\

\citep{gaebele_tpl_2025} & 2-PA & BEM & - & 2 & PS & PI & - & DYN, EFF & - & PTO, GEO & IRR
& LOC & ECON & PTO & PTO, CTL, DIM & 5 \\

\citep{devin_high-dimensional_2024} & 1-PA & FIT & - & 1 & PS & UNS & - & DYN, EFF & - & - & REG
& DOE, BF, LOC & PWR & PTO & PTO, CTL & 5 \\

\citep{grasberger_control_2024} & OSW & BEM & - & 1 & PS & PI & - & DYN & - & GEO & IRR
& BF, LOC & ECON & AMP, PTO & DIM, PTO, CTL & 2 \\

\citep{herber_dynamic_2014} & 1-PA & FIT & LD & 1 & T & UNS & - & - & - & - & IRR& BF, LOC & PWR & AMP, PTO & DIM, CTL & 2 \\

\citep{lin_fast_2025} & 1-PA & BEM & LD & 1 & F+ & UNS & - & - & - & - & IRR+
& BF, PMP & ECON & AMP, PTO & DIM & 1 \\

\citep{abdulkadir_optimal_2024} & 1-PA array & BEM & - & 3 & T & UNS & - & - & - & - & IRR
& PMP & PWR & PTO & CTL & - \\

%\citep{zhang_performance_2024} & 3-PA & MEEM & LD & 3 & F & P & DYN & - & - & - & IRR
%& BF & PWR & - & CTL & 2 \\

%\citep{wen_shape_2018} & 1-PA & BEM & - & 1 & F & P & - & - & - & - & IRR+
%& DOE, LOC & PWR & STA & DIM, SHP & 3 \\

\hline 
%\end{tabular}
\caption{Comparison of previous WEC optimization studies' disciplinary scope, modeling fidelity, and optimization formulation.
See \Cref{remdo:sec:appendix-lit-legend} for legend definitions.}
\label{remdo:tab:lit}
%\fillandplacepagenumber
\endlongtable
}
\endgroup
%\end{landscape}

% Prior WEC optimization across discipline
\ifdefined\DISSERTATION
    A handful of studies perform WEC optimization spanning several disciplines.
    \citet{gaudin_single_2021} optimizes an array of 20 submerged cylinders, holding the device design constant while varying site, layout, and mooring with 120 degrees of freedom.
    \citet{khanal_multi-objective_2024} optimizes a four-WEC array's layout, hydrodynamics, controls, and economics (14 design variables).
    \citet{edwards_optimisation_2022} performs point-absorber shape optimization with stability constraints, minimizing surface area subject to power-extraction targets at radiation, amplitude, and steepness limits.
    \citet{garcia-teruel_reliability-based_2021} and the smaller-scale \citet{garcia-teruel_design_2022} apply multi-objective optimization to trade off power with fatigue load and device volume respectively.
    \citet{cotten_multi-objective_2022} performs a similar fatigue-load study for an attenuator-style WEC with eight design variables.
    \citet{abdulkadir_control_2024} performs an eight-variable shape optimization using semi-analytical hydrodynamics (MEEM) but no other disciplines beyond hydrodynamics and dynamics.

    All seven of these large-scale studies \citep{khanal_multi-objective_2024,gaudin_single_2021,edwards_optimisation_2022,garcia-teruel_reliability-based_2021,garcia-teruel_design_2022,cotten_multi-objective_2022,abdulkadir_control_2024} use the genetic algorithm, a population-based heuristic.
    The alternative is local optimization, which lacks heuristic exploration but leverages smoothness and scales well to many design variables \citep{martins_engineering_2022}.
    \citet{housner_numerical_2024} attempted local optimization for a terminator-shape study and could only handle 3 of an intended 13 design variables due to time-domain simulation cost;
    \citet{garcia-teruel_reliability-based_2021,garcia-teruel_design_2022,cotten_multi-objective_2022} avoided that issue by simulating primarily in the frequency domain, falling back to time domain only for PTO constraints.
    Design-of-experiment alternatives such as the Taguchi method \citep{al_shami_parameter_2019} enable systematic exploration without traditional optimizers.

    Structural optimization of WECs faces an analogous compute bottleneck: FEA-based optimization is avoided in favor of commercial decoupled tools \citep{an_optimal_2024}, brute-force sweeps over analytical stress \citep{ambuhl_reliability-based_2014}, load-magnitude proxies in place of stress \citep{nguyen_theoretical_2024,ferri_balancing_2014,garcia-teruel_reliability-based_2021,cotten_multi-objective_2022}, or manual iteration \citep{RM3,mi_multi-scale_2025}.
    The last two studies \citep{RM3,mi_multi-scale_2025} notably encompass mooring, powertrain, controls, structures, and detailed economics.

    The control co-design (CCD) literature \citep{garcia-sanz_control_2019,coe_useful_2023} emphasizes drivetrain dynamics and generator efficiency due to their strong effect on electrical power.
    Recent joint PTO-and-controller optimizations \citep{rosati_control_2023,son_performance_2016,anderson_re-imagining_2024,devin_high-dimensional_2024,grasberger_control_2024,gaebele_tpl_2025} often leverage the pseudo-spectral method to embed time-domain constraints in frequency-domain hydrodynamics.
    Of these, \citet{gaebele_tpl_2025} directly addresses RM3, finding that a 0.55 scaling factor optimizes power per unit surface area with frequently active continuous-torque constraints.
    \citet{lin_fast_2025} introduces a CCD formulation faster than PS by leveraging Pontryagin's maximum principle analytically; \citet{abdulkadir_optimal_2024} use PMP for PTO-constrained array dynamics; \citet{herber_dynamic_2014} performs direct-transcription co-optimization of cylinder draft and radius with control.

    Several studies have independently concluded that significantly undersizing the PTO and accepting force saturation is favorable.
    \citet{mcgilton_optimal_2024} reports only 3\% energy loss from a 70\% force-rating reduction across multiple device archetypes (RM3, RM5), sites, and scales, supporting analytical predictions \citep{devin_high-dimensional_2024}.
    \citet{coe_maybe_2021} argues both for a smaller PTO and a smaller physical device, reasoning that ``due first to the high variability and strong positive skew of ocean wave power and second to the filtering behavior of a WEC, a smaller WEC may actually have better economics in terms of average annual energy, capacity factor, and costs.''
    None of these studies perform full economic analysis to confirm that the cost savings outweigh the power decrease --- a gap this paper addresses.
    \else
    Of the multi-discipline studies in \Cref{remdo:tab:lit}, the largest in scope are those of \citet{khanal_multi-objective_2024,gaudin_single_2021,edwards_optimisation_2022,garcia-teruel_reliability-based_2021,garcia-teruel_design_2022,cotten_multi-objective_2022}, spanning four-WEC array layout, hydrodynamics, controls, economics, high-dimensional site/layout/mooring, shape with stability constraints, and power-versus-structure trade-offs.
    All large-scale ($\geq 8$ design variables) WEC optimizations use the genetic algorithm; local optimization is limited to small design spaces or frequency-domain settings due to time-domain simulation cost \citep{housner_numerical_2024,garcia-teruel_reliability-based_2021,garcia-teruel_design_2022,cotten_multi-objective_2022,martins_engineering_2022}.
    Structural optimization faces a compute bottleneck, with FEA-based optimization usually replaced by manual iteration \citep{RM3,mi_multi-scale_2025}, commercial decoupled tools \citep{an_optimal_2024}, brute-force sweeps over analytical stress \citep{ambuhl_reliability-based_2014}, or load-magnitude proxies \citep{nguyen_theoretical_2024,ferri_balancing_2014}.
    The control co-design literature \citet{garcia-sanz_control_2019,coe_useful_2023} has produced joint PTO-and-controller optimizations \citet{rosati_control_2023,son_performance_2016,anderson_re-imagining_2024,devin_high-dimensional_2024,grasberger_control_2024,gaebele_tpl_2025} often leveraging the pseudo-spectral method;
    \citet{gaebele_tpl_2025} directly optimizes RM3 and finds a 0.55 dimensional scaling factor optimal.
    Several studies independently conclude that undersizing the PTO and accepting force saturation is favorable: \citet{mcgilton_optimal_2024} reports only 3\% energy loss from a 70\% force-rating reduction across archetypes, sites, and scales, supporting analytical predictions made in other studies \citep{coe_maybe_2021,devin_high-dimensional_2024}.
    None perform a full economic analysis to confirm cost savings exceed power loss, a gap this paper addresses.
    \fi

    % Gaps
    \ifdefined\DISSERTATION
    The literature review reveals that no existing optimization study includes dimensional, structural, and PTO design variables simultaneously.
    The closest are the studies by \citet{ambuhl_reliability-based_2014} (dimensional and structural) and by \citet{grasberger_control_2024} (dimensional and PTO).
    There is also a gap of work that uses a scalable optimization framework to trade off power production with survivability while accurately modeling stress as a function of bulk dimensions, rather than using force as a proxy.
    The closest is \citet{rosati_control_2023}, which minimizes a cost proxy incorporating PTO and valve material cost but assumes constant structural material cost.
    Given the coupling among these disciplines, a system-level optimization considering them all is needed.
    Part of the issue is the lack of a multidisciplinary WEC simulation platform sufficiently fast for optimization: the popular time-domain hydrodynamics tool WEC-Sim \citep{ruehl_wec-simwec-sim_2024} is too slow for large-scale optimization \citep{housner_numerical_2024}, while the newer pseudo-spectral tool WecOptTool \citep{coe_initial_2020} is faster and suited to control co-design but currently limited to hydrodynamics and control.
    Finally, several studies note that cost-model uncertainty may significantly affect results, evidencing the need for sensitivity analysis.
\else
    No existing optimization study simultaneously includes dimensional, structural, and PTO design variables: the closest are the studies by \citet{ambuhl_reliability-based_2014,grasberger_control_2024}, covering dimensional-and-structural and dimensional-and-PTO formulations respectively.
    A parallel gap exists in trading off power against survivability with stress modeled as a function of bulk dimensions rather than force as a proxy.
    The lack of a multidisciplinary WEC simulation platform fast enough for optimization is part of the issue (e.g., WEC-Sim \citep{ruehl_wec-simwec-sim_2024} is too slow at scale \citep{housner_numerical_2024}, and WecOptTool \citep{coe_initial_2020} is currently limited to hydrodynamics and control).
    Cost-model uncertainty in prior work \citep{mccabe_multidisciplinary_2022} evidences the need for sensitivity analysis.
\fi

%%%%%%%%%%%%%%%%%%%%%%%%%%%%%%%%%%%%%%%%%%%%%%%%%%
\subsection{Relevance of Multidisciplinary Design Optimization}
%(Why MDO):
\ifdefined\DISSERTATION
    Multidisciplinary Design Optimization (MDO) is a quantitative methodology for formulating, solving, and interpreting complex engineering design problems involving multiple coupled disciplines \citep{agte_mdo_2010}.
    The method emphasizes that comprehensive consideration of inter-disciplinary coupling typically yields designs that outperform those found by sequential single-discipline optimization \citep{martins_multidisciplinary_2013}.
    MDO falls under the broader umbrella of concurrent design and, when controls is one of the disciplines, constitutes a form of control co-design.
    The field originated in structural optimization, is well-established in aerospace, and is gaining traction in the broader energy industry, particularly in floating offshore wind \citep{abbas_control_2024,jasa_effectively_2022,patryniak_multidisciplinary_2022}.

    Despite recognized need and growing interest in concurrent design and control co-design for WECs \citep{mi_multi-scale_2025,ringwood_empowering_2023}, MDO has not yet caught on in the wave energy space, preventing potential techno-economic breakthroughs.
    Only two existing studies \citep{mccabe_multidisciplinary_2022,khanal_multi-objective_2024} approach WEC design explicitly from an MDO perspective.
    The present study builds directly on \citet{mccabe_multidisciplinary_2022}, a preliminary MDO formulation by the same authors that used substantially less accurate models and fewer design variables.
    \citet{khanal_multi-objective_2024} adds a Sobol global sensitivity analysis to a basic MDO formulation.
    Neither provides a comprehensive description of how the MDO approach influences model development and problem formulation, which this paper provides alongside results and insights from a much more comprehensive WEC optimization.

    The remaining WEC papers discussed as multidisciplinary do not explicitly invoke MDO and sometimes do not adhere to the principle of adequately capturing inter-disciplinary coupling.
    For instance, \citet{edwards_optimisation_2022} minimizes surface material for only the subset of designs achieving maximum power, which fails to account for the power--cost tradeoff: the minimum cost-to-power ratio may occur at a design that sacrifices some power for significant cost savings.

    Semi-analytical models offer a uniquely suitable combination of speed and accuracy for early-stage WEC optimization, where the model must be run thousands of times and where smooth dependence on inputs facilitates accurate gradients \citep{mccabe_development_2026}.
    This study uses the MDOcean framework \citep{mccabe_development_2026}, which integrates semi-analytical models across hydrodynamics, dynamics and control, structures, and economics, to enable a WEC optimization of unprecedented scope.
\else
    Multidisciplinary Design Optimization (MDO) is a quantitative methodology for formulating, solving, and interpreting design problems involving multiple coupled disciplines \citep{agte_mdo_2010}, highlighting that comprehensive consideration of coupling typically outperforms sequential single-discipline optimization \citep{martins_multidisciplinary_2013}.
    Originating in structural optimization and established in aerospace, MDO is gaining traction in floating offshore wind \citep{abbas_control_2024,jasa_effectively_2022,patryniak_multidisciplinary_2022} but has not yet caught on in wave energy despite recognized need \citep{mi_multi-scale_2025,ringwood_empowering_2023}.
    Only two existing studies approach WEC design explicitly from an MDO perspective: \citet{mccabe_multidisciplinary_2022}, a preliminary formulation by the present authors using less accurate models and fewer design variables, and \citet{khanal_multi-objective_2024}, which adds Sobol global sensitivity.
    Other multidisciplinary WEC papers do not explicitly invoke MDO and sometimes neglect the coupling principle: for example, \citet{edwards_optimisation_2022} minimizes surface material only over designs achieving maximum power, missing the power--cost tradeoff where the minimum cost-to-power ratio may sacrifice some power for significant cost savings.

    This study uses the MDOcean framework \citep{mccabe_development_2026}, which integrates semi-analytical models across hydrodynamics, dynamics and control, structures, and economics, to enable a WEC optimization of unprecedented scope.
\fi

%%%%%%%%%%%%%%%%%%%%%%%%%%%%%%%%%%%%%%%%%%%%%%%%%%
\subsection{Paper Contributions and Roadmap}
\ifdefined\DISSERTATION
    This paper applies the MDOcean framework \citep{mccabe_development_2026} to what is believed to be the most comprehensive WEC techno-economic design optimization to date.
    The contributions are: a multidisciplinary problem formulation simultaneously varying 12 design variables across bulk dimensions, structural thicknesses, and powertrain ratings, subject to static, dynamic, structural, and operational constraints;
    a solution methodology combining sequential quadratic programming with multistart for single-objective optimization and the epsilon-constraint method seeded with pattern search for multi-objective optimization;
    single-objective optimization results with local and global sensitivity analysis identifying drivers of LCOE, cost, and power;
    multi-objective optimization results revealing a non-convex Pareto front with a cusp corresponding to a force-saturation to power-saturation strategy shift, and quantitative design heuristics across the front;
    and design recommendations distinguishing utility-scale from Powering-the-Blue-Economy applications.
    The framework is scalable to additional disciplines and design variables and is ultimately intended to enable systematic comparison across WEC architectures.

    The remainder of the paper is organized as follows.
    \Cref{remdo:sec:optim-methods} describes the optimization methodology including the underlying model topology, although modeling details are published separately \citep{mccabe_development_2026}.
    \Cref{remdo:sec:formulation} develops the optimization problem formulation.
    \Cref{remdo:sec:optim-process} details the solution procedure.
    \Cref{remdo:sec:results} presents single-objective optimization, sensitivity analysis, and multi-objective optimization results,
    then interprets the trends, presents design takeaways, and outlines future work.
\else
    This paper applies the MDOcean framework \citep{mccabe_development_2026} to what is believed to be the most comprehensive WEC techno-economic design optimization to date.
    The contributions are: a multidisciplinary problem formulation simultaneously varying 12 design variables across bulk dimensions, structural thicknesses, and powertrain ratings subject to static, dynamic, structural, and operational constraints;
    a solution methodology combining sequential quadratic programming with multistart and an epsilon-constraint-plus-pattern-search procedure for multi-objective optimization;
    single- and multi-objective results with sensitivity analysis revealing drivers of LCOE, cost, and power;and design recommendations distinguishing utility-scale from \emph{Powering the Blue Economy} applications.

    The remainder of the paper is organized as follows.
    \Cref{remdo:sec:optim-methods} describes the optimization methodology, and \citet{mccabe_development_2026} describes the underlying simulation model.
    \Cref{remdo:sec:formulation} develops the problem formulation and \Cref{remdo:sec:optim-process} the solution procedure.
    \Cref{remdo:sec:results} presents and interprets results and outlines future work.
\fi
